\section{实验目的}

推荐系统是解决信息过载问题的核心技术之一，其核心任务是根据用户的历史行为预测其对未交互物品的偏好。本实验以电影评分数据为载体，围绕协同过滤的两大范式——基于模型的方法（矩阵分解）与基于记忆的方法（近邻法）——展开系统性的实现与对比，具体目标如下：

\begin{enumerate}
    \item \textbf{掌握协同过滤的核心算法}：从原理层面理解矩阵分解（FunkSVD、SVD++）与基于邻域的协同过滤（Item-KNN）的数学推导与优化方法，并通过代码实现加深认知。
    \item \textbf{体会稀疏性对算法的影响}：在评分密度极低（约1.5\%）的真实场景下训练模型，观察数据稀疏性如何制约不同范式的算法表现。
    \item \textbf{建立算法对比的分析框架}：从预测精度（RMSE）、训练效率（时间）、空间开销三个维度系统对比不同算法，分析各方法在不同条件下的适用边界。
    \item \textbf{理解超参数的作用机制}：通过网格搜索探索隐因子维度、学习率、正则化强度、近邻数等超参数对模型性能的影响，理解过拟合-欠拟合的权衡。
\end{enumerate}
\section{数据集分析}

\subsection{数据格式}

数据集包含三个文件：\texttt{train.txt}（训练集）、\texttt{test.txt}（测试集）和 \texttt{ResultForm.txt}（结果模板）。数据格式以用户为单位组织：每个用户条目以 \texttt{<user id>|<评分数>} 开头，随后逐行列出该用户评价过的 \texttt{<item id> <score>}（训练集）或 \texttt{<item id>}（测试集，无评分）。

\subsection{基本统计量}

表~\ref{tab:stats} 汇总了数据集的核心统计指标。

\begin{table}[htbp]
    \centering
    \caption{数据集基本统计量}
    \label{tab:stats}
    \begin{tabular}{@{}lr@{}}
        \toprule
        \textbf{指标} & \textbf{数值} \\
        \midrule
        训练集用户数     & 598 \\
        训练集物品数     & 9,077 \\
        训练集评分数     & 90,854 \\
        测试集 (用户, 物品) 对数 & 9,982 \\
        测试集用户数     & 610 \\
        测试集物品数     & 3,618 \\
        全部用户数       & 610 \\
        全部物品数       & 9,724 \\
        可能的 (用户, 物品) 对数 & 5,931,640 \\
        \textbf{稀疏度}  & \textbf{98.47\%} \\
        \bottomrule
    \end{tabular}
\end{table}

数据集极其稀疏：已知评分仅占全部可能评分对的 1.53\%，这对算法的泛化能力提出了较高要求。值得注意的是，测试集包含 9,982 个待预测对，其中存在 12 个全新用户和 647 个全新物品，构成一定程度的冷启动挑战。

\subsection{评分分布}

训练集中评分范围为 $[10, 100]$，评分统计量如表~\ref{tab:rating_stats} 所示，每 10 分一档的分布见表~\ref{tab:rating_dist}。

\begin{table}[htbp]
    \centering
    \caption{评分统计量}
    \label{tab:rating_stats}
    \begin{tabular}{@{}lcccccc@{}}
        \toprule
         & 均值 & 方差 & 标准差 & 最小值 & 最大值 & 中位数 \\
        \midrule
        评分 & 69.88 & 431.81 & 20.78 & 10 & 100 & 70 \\
        \bottomrule
    \end{tabular}
\end{table}

\begin{table}[htbp]
    \centering
    \caption{评分分布（10 分档）}
    \label{tab:rating_dist}
    \begin{tabular}{@{}lrrlrr@{}}
        \toprule
        \textbf{分数} & \textbf{数量} & \textbf{占比} & \textbf{分数} & \textbf{数量} & \textbf{占比} \\
        \midrule
        10 & 1,221 & 1.34\% & 60 & 18,119 & 19.94\% \\
        20 & 2,552 & 2.81\% & 70 & 11,996 & 13.20\% \\
        30 & 1,603 & 1.76\% & 80 & 24,177 & 26.61\% \\
        40 & 6,852 & 7.54\% & 90 & 7,742  & 8.52\%  \\
        50 & 5,067 & 5.58\% & 100 & 11,525 & 12.69\% \\
        \bottomrule
    \end{tabular}
\end{table}

评分分布呈明显的\textbf{高分偏态}：80 分是众数（26.61\%），60--100 分区间合计占 81.0\%，中位数 70 分，而 50 分以下仅占 19.0\%。这符合电影评分场景的常见模式——用户倾向于观看并高分评价自己喜欢的电影，低分评价相对稀缺。

\subsection{用户与物品活跃度}

\begin{table}[htbp]
    \centering
    \caption{用户与物品的评分数分布}
    \label{tab:activity}
    \begin{tabular}{@{}lcccc@{}}
        \toprule
        \textbf{统计量} & \textbf{均值} & \textbf{最少} & \textbf{最多} & \textbf{中位数} \\
        \midrule
        用户评分数   & 151.93 & 15   & 2,678 & 59 \\
        物品被评数   & 10.01  & 1    & 280   & 3  \\
        \bottomrule
    \end{tabular}
\end{table}

用户活跃度分布不均：人均评价 152 部电影，但中位数仅 59 部，说明少数高活跃用户贡献了大量评分。物品端呈现更极端的\textbf{长尾分布}：65.6\% 的物品仅获得 5 次及以下评价（5,959 个），中位数仅为 3 次，而头部物品最多被评价 280 次。这种长尾特性是邻域方法（Item-KNN）表现不佳的根本原因——大量冷门物品与其他物品的共现次数过少，相似度估计不可靠。

\subsection{冷启动分析}

测试集中有 12 个用户（占 2.0\%）和 647 个物品（占 17.9\%）从未在训练集中出现。对于新用户，模型只能依赖全局均值进行预测；对于新物品，基于物品相似度的 KNN 方法完全失效，而矩阵分解方法可通过偏置项给出合理的基线估计。
\section{算法描述}
\section{实验设计}
\section{实验结果}
\section{分析与讨论}